% \iffalse
%
\ProvidesFile{tree-sitter.dtx}
%
%\ProvidesPackage{tree-sitter}
[2026/08/04 v1.0.0 LaTex bindings for tree-sitter]
%
%<*driver>
\documentclass{ltxdoc}
\EnableCrossrefs
\CodelineIndex
\begin{document}
  \DocInput{tree-sitter.dtx}
\end{document}
%</driver>
% \fi
%
% \GetFileInfo{tree-sitter.dtx}
% \title{LaTex bindings for tree-sitter}
% \date{\filedate\qquad\fileversion}
% \maketitle
% \begin{abstract}
% This is the implementation of the simple-tree example.
% \end{abstract}
%
%
% \tableofcontents
%
% \section{Implementation}
%
%    \begin{macrocode}
\RequirePackage{shellesc}
\RequirePackage{pgfkeys}
\RequirePackage{xcolor}
\RequirePackage{fancyvrb}

\pgfkeys{
  /tree-sitter/.is family,
  /tree-sitter,
  /tree-sitter/language/.store in=\TreeSitterLanguage,
  /tree-sitter/language/.initial=c,
}

\newcommand{\treeSitterLanguage}[1]{\pgfkeys{/tree-sitter/language=#1}}

\newcommand{\buildScopeTreeSitter}[1]{source.#1}

\makeatletter
\newcounter{TreeSitter@counter}
\newcommand{\TreeSitterCounter}{\arabic{TreeSitter@counter}}
\makeatother

\newenvironment{tree-sitter}[1][]
{
  \pgfkeys{/tree-sitter, #1}
  \stepcounter{TreeSitter@counter}
  \ifx\ShellEscape\undefined
    \PackageError{tree-sitter}{Shell escape is not enabled}
    {add shell-escape for build tex file}
  \fi
  \ifnum\ShellEscapeStatus=0
    \PackageError{tree-sitter}{Shell escape is disabled}
    {add shell-escape for build tex file}
  \fi
  \VerbatimEnvironment
  \begin{VerbatimOut}{\jobname-tree-sitter-input-\TreeSitterCounter.}%
}{
  \end{VerbatimOut}
  \ShellEscape{cat "\jobname-tree-sitter-input-\TreeSitterCounter."
  | tree-sitter-highlight highlight -L --scope "\buildScopeTreeSitter{\TreeSitterLanguage}"
  > \jobname-tree-sitter-\TreeSitterCounter.}
  \input{\jobname-tree-sitter-\TreeSitterCounter.}
  \ShellEscape{rm \jobname-tree-sitter-input-\TreeSitterCounter.}
  \ShellEscape{rm \jobname-tree-sitter-\TreeSitterCounter.}
}

\newcommand{\treesitterinput}[2][]{%
  \stepcounter{TreeSitter@counter}%
  \ifx\ShellEscape\undefined
    \PackageError{tree-sitter}{Shell escape is not enabled}
  \fi
  \ifnum\ShellEscapeStatus=0
    \PackageError{tree-sitter}{Shell escape is disabled}
  \fi
  \newcommand{\inputfilename}{#2}%
  \ShellEscape{tree-sitter-highlight highlight --latex "\inputfilename"
  > \jobname-tree-sitter-\TreeSitterCounter.}%
  \input{\jobname-tree-sitter-\TreeSitterCounter.}%
  \ShellEscape{rm \jobname-tree-sitter-\TreeSitterCounter.}%
}

%    \end{macrocode}
%
% \Finale
